\documentclass{situationist-preprint}

\title{The Footbridge Case}
\msid{ETH-03-b}
\msfamily{ETH-03}
\msdiscipline{Normative and Applied Ethics}
\mstype{Companion manuscript}
\mskeywords{trolley problem; Footbridge; redirected threat; Loop; numbers; using a person}

\begin{document}
\maketitle

\begin{abstract}
ETH-03-a takes the impermissibility of pushing the man in Footbridge as one of two data and proves that the doctrine of the redirected threat is the only rule that fits the data and two further adequacy conditions. We examine the Footbridge verdict itself. We show that, under the doctrine, it does not depend on the man's body being used as a means, which Loop shares, nor on the number of people saved, and we defend the second consequence, which is the less familiar one.
\end{abstract}

\section{What makes Footbridge impermissible}

In Footbridge \citep{Thomson1985}, a bystander can push a large man from a bridge into the path of a trolley; his body will stop it, killing him and saving five. By Definition~2.2 of ETH-03-a, the doctrine of the redirected threat (DRT) permits an act in a trolley case if and only if it redirects an existing threat and saves more than it kills. Pushing the man saves more than it kills, but it does not redirect the trolley: the trolley continues on its course, and the man, who was not in its path, is put there. The act introduces a new threat to him, and the DRT forbids it. This is the verdict Principle~P1 of ETH-03-a requires, and Proposition~3.3 of ETH-03-a records that the DRT delivers it.

The usual explanation is different: the man is used as a means, since it is the impact of his body that stops the trolley. Loop shows that this is not what matters. In Loop, the one on the side track also stops the trolley with his body, and his being struck is equally a means to the survival of the five; yet ETH-03-c proves (Theorem~3.1 of ETH-03-c) that diverting the trolley in Loop is permissible, from the principle that redirecting an existing threat so that it kills fewer is permissible (Principle~P1 of ETH-03-c) and the irrelevance of track geometry beyond the point at which the redirected threat reaches its victims (Principle~P2 of ETH-03-c).

Footbridge and Loop agree in the use of a body as a means and differ in one respect: where the victim was before the act. The one in Loop is already on a track to which the trolley can be turned; the act changes the trolley's path, and he is in the new path. The man in Footbridge is on no track; the act changes his position, not the trolley's. This difference is not one of track geometry beyond the victim, so Principle~P2 of ETH-03-c does not apply to it, and it is exactly the difference that the first feature of Definition~2.1 of ETH-03-a records. The turntable case of ETH-03-a, Section~5, in which the agent rotates a platform carrying the one into the trolley's path, falls on the Footbridge side for the same reason.

\section{Numbers do not rescue the push}

By Principle~P3 of ETH-03-a, the verdict on a trolley case depends only on whether the act redirects an existing threat and whether it saves more than it kills. A variant of Footbridge in which pushing the man saves fifty, or five thousand, agrees with Footbridge on both features, and by Theorem~3.1 of ETH-03-a every adequate rule forbids it.

\begin{objection}
This is incredible. If pushing one man would save a city, it would be permissible.
\end{objection}

\begin{reply}
The objection appeals to a threshold above which constraints give way. Such a threshold would make the verdict depend on the exact number saved, which Principle~P3 excludes, and it would leave the victim's claim to depend on how many others happen to be in danger, which is precisely what the Footbridge verdict denies. The man on the bridge has the same claim not to be made into a shield whether five or fifty are threatened. Cases in which the stakes are a city are, moreover, rarely trolley cases in the sense of Definition~2.1 of ETH-03-a, since they involve uncertainty, institutions and many agents; the DRT is silent about them, as ETH-03-a notes for triage. Within its class of cases, the doctrine is absolute, and we accept that consequence.
\end{reply}

\section{Conclusion}

Pushing the man is impermissible because it introduces a threat to someone who was not under it, not because it uses him, and the verdict holds however many are saved. Footbridge and Loop, which agree in using a person's body, differ in exactly the feature the doctrine of the redirected threat makes decisive.

\begin{thebibliography}{9}\small
\bibitem[Thomson(1985)]{Thomson1985} Thomson, J.~J. (1985). The trolley problem. \emph{The Yale Law Journal}, 94(6), 1395--1415.
\end{thebibliography}

\end{document}
